Let \(C_Q\), \(D\), and \(R\) denote, respectively, the covariate block, binary treatment, and outcome in \(\mathrm{lem:qiu\text{-}external\text{-}ate\text{-}characterization}\).  Relabel the variables for the purpose of comparing signatures by
\[
T_0=C_Q,\qquad X=\{D\},\qquad Y=\{R\},\qquad
V=C_Q\mathbin{\dot\cup}\{D,R\}.
\tag{1}
\]
By \(\mathrm{def:observed\text{-}data\text{-}map}\), the two entries of the present observation map under (1) are
\[
\operatorname{Obs}_{T_0}(M)
=\bigl(P_M(V\mid S=1),P_M(T_0)\bigr)
=\bigl(P_M(C_Q,D,R\mid S=1),P_M(C_Q)\bigr).
\tag{2}
\]
The cited lemma states that Qiu et al. use a selected observational law together with an external unbiased covariate margin.  Thus (2) proves equality of the two data-slot patterns after relabeling.  It does not assert equality, inclusion, or even nonempty intersection of the two model classes.

We next verify the estimand map.  Suppose only for this calculation that \(R\) has a finite numeric state space and \(D\) is binary.  Every sum below is finite, so no integrability or limiting condition is needed.  The definitions of \(\theta_{\mathrm{ATE}}\) and \(\psi_{x,y}\) give, equation by equation,
\[
\begin{aligned}
\theta_{\mathrm{ATE}}(M)
&=E_M[R\mid\operatorname{do}(D=1)]
  -E_M[R\mid\operatorname{do}(D=0)]\\
&=\sum_{r\in\mathcal X_R}r\,
  P_M(R=r\mid\operatorname{do}(D=1))
  -\sum_{r\in\mathcal X_R}r\,
  P_M(R=r\mid\operatorname{do}(D=0))\\
&=\sum_{r\in\mathcal X_R}r\,\psi_{1,r}(M)
  -\sum_{r\in\mathcal X_R}r\,\psi_{0,r}(M).
\end{aligned}
\tag{3}
\]
Hence the finite table \(\boldsymbol\psi(M)\), whose coordinatewise role is made explicit in \(\mathrm{def:full\text{-}kernel\text{-}collection}\), determines \(\theta_{\mathrm{ATE}}(M)\) through the displayed linear map.

The converse to this last implication fails inside a strictly positive finite Markovian DAG subclass.  Consider the graph \(D\to R\) with \(S\) isolated, take \(X=\{D\}\), \(Y=\{R\}\), and let \(T_0=\varnothing\).  Let \(U_D\) and \(U_S\) be independent Bernoulli\((1/2)\) variables, and let \(U_R\) be independent of them and uniform on \(\{1,\ldots,12\}\).  In model \(M_j\), set
\[
D=U_D,\qquad S=U_S,
\qquad
R_1=\mathbf 1\{U_R\le 3+6D\},
\qquad
R_2=\mathbf 1\{U_R\le 4+6D\}.
\tag{4}
\]
Each model is finite and recursive.  Its exogenous variables are mutually independent, so it is Markovian and therefore semi-Markovian; its structural equations follow the displayed DAG.  Also \(P_{M_j}(S=1)=1/2>0\).  Independence of \(S\) from \((D,R_j)\) yields the selected cell tables
\[
\begin{array}{c|cccc}
&(D,R_j)=(0,1)&(0,0)&(1,1)&(1,0)\\ \hline
j=1&1/8&3/8&3/8&1/8\\
j=2&1/6&1/3&5/12&1/12.
\end{array}
\tag{5}
\]
Every cell in (5) is positive; in fact both models satisfy the selected-cell lower bound \(1/12\).  Thus the positivity assertion uses no boundary point.  Intervention replaces the equation \(D=U_D\) and leaves the equation for \(R_j\) unchanged, so (4) gives
\[
\begin{array}{c|cc|c}
&P_{M_j}(R=1\mid\operatorname{do}(D=0))
&P_{M_j}(R=1\mid\operatorname{do}(D=1))
&\theta_{\mathrm{ATE}}(M_j)\\ \hline
j=1&3/12&9/12&6/12\\
j=2&4/12&10/12&6/12.
\end{array}
\tag{6}
\]
Both scalar effects equal \(1/2\), while \(\psi_{0,1}(M_1)=1/4\ne1/3=\psi_{0,1}(M_2)\).  Therefore the linear map in (3) is not injective even on the stated positive finite DAG subclass: identification of the scalar effect alone cannot decide every coordinate of \(\boldsymbol\psi\).

It remains to establish the two scope separations.  By \(\mathrm{lem:qiu\text{-}external\text{-}ate\text{-}characterization}\), Qiu et al.'s theorem assumes unconfoundedness, treatment overlap on the covariate support, and a pairwise distinguishability condition on specified classes of propensity, covariate--potential-outcome, and selection laws.  Its conclusion identifies the single scalar \(\theta_{\mathrm{ATE}}\), not every cell \(\psi_{x,y}\) for arbitrary semi-Markovian \(\mathcal G\).  Equation (6) additionally shows why the scalar conclusion cannot be inverted into the full-table conclusion.  Thus the cited theorem does not subsume the present arbitrary-cell graph problem.

Conversely, \(\mathrm{def:model\text{-}class}\) restricts the present class to finite-state recursive semi-Markovian models and imposes only the listed graph, selection, and selected-cell positivity conditions.  The cited lemma states both that Qiu et al.'s exposition includes continuous covariate--outcome laws with densities and that its identification criterion is indexed by separately specified probability classes.  Continuous-law examples are not members of the present finite-state class.  Moreover, a graph \(\mathcal G\) and the membership conditions in \(\mathrm{def:model\text{-}class}\) do not supply a chosen triple of probability classes or assert Qiu et al.'s pairwise distinguishability condition.  Hence the present graph class does not subsume the cited setup.  Together with (2), (3), and (6), these observations prove precisely an input-shape and estimand-map analogy, the failure of the reverse estimand map, and non-subsumption in both directions, without claiming a literal model-class intersection.